\section{改进二：非线性为什么会导向 MLP？}

之前每条记忆的贡献是：
\[
\vect{v}_i\times\underbrace{(\vect{k}_i^\top\vect{x})}_{\text{原始相似度}}.
\]
例如正确 Key 的相似度是 $1$，无关 Key 是 $0.6$。线性加权得到 $1$ 和 $0.6$，错误记忆仍贡献很多。

最简单的办法：\textbf{把相似度平方}！
\[
1^2=1,\qquad 0.6^2=0.36.
\]
如果立方：
\[
1^3=1,\qquad 0.6^3=0.216.
\]
正确记忆保持为 $1$，干扰越来越小。

\begin{center}
\begin{tikzpicture}
\begin{axis}[
 width=.89\linewidth,height=5.4cm,
 axis lines=left,xlabel={相似度 $s$},ylabel={激活权重},
 xmin=0,xmax=1.02,ymin=0,ymax=1.06,
 xtick={0,.2,.4,.6,.8,1},ytick={0,.2,.4,.6,.8,1},
 legend style={at={(.02,.98)},anchor=north west,font=\small,draw=none,fill=white},
 grid=major,grid style={linegray!55}, tick label style={font=\small},
 ]
 \addplot[blue,very thick,domain=0:1,samples=100] {x};
 \addlegendentry{线性 $s$}
 \addplot[orange,very thick,domain=0:1,samples=100] {x^2};
 \addlegendentry{平方 $s^2$}
 \addplot[teal,very thick,domain=0:1,samples=100] {x^3};
 \addlegendentry{立方 $s^3$}
 \addplot[dashed,gray] coordinates {(.6,0)(.6,1)};
\end{axis}
\end{tikzpicture}
\end{center}

为了把负相似度也压成零，可以使用：
\begin{equation}
\sigma(s)=\relu(s)^2=\bigl(\max(0,s)\bigr)^2.
\end{equation}
我们于是得到新的记忆检索器：
\begin{equation}
\boxed{f(\vect{x})=\sum_i\vect{v}_i\,\sigma(\vect{k}_i^\top\vect{x}).}
\end{equation}
将所有 Key 堆成矩阵 $K$，所有 Value 堆成矩阵 $V$，就可以简写为：
\begin{equation}
\boxed{f(\vect{x})=V^\top\sigma(K\vect{x}).}
\end{equation}
这里约定 $K$ 的每一行是一条 Key，$V$ 的每一行是一条 Value；激活函数 $\sigma$ 对向量中的各元素分别作用。

\begin{center}
\begin{tikzpicture}[>=Latex,node distance=0.55cm,n/.style={draw=linegray,rounded corners=3pt,fill=paper,minimum width=92mm,minimum height=11mm,align=center,font=\small}]
\node[n] (q) {输入 Query $\vect{x}$};
\node[n,below=of q] (s) {第一层：$K\vect{x}$，计算每条记忆的匹配分数};
\node[n,below=of s] (a) {非线性：$\sigma(K\vect{x})$，压低不相关记忆的激活};
\node[n,below=of a] (v) {第二层：$V^\top\sigma(K\vect{x})$，组合 Value 得到输出};
\draw[->,blue,thick] (q)--(s);\draw[->,blue,thick] (s)--(a);\draw[->,blue,thick] (a)--(v);
\end{tikzpicture}
\end{center}

这恰好是一个\textbf{两层 MLP} 的典型数学形式：一次线性变换、一次非线性激活、再一次线性变换。

\begin{conclusion}
从记忆角度看，MLP 的第一层可以被解释为\textbf{匹配模式}，激活函数负责\textbf{筛选}，第二层负责\textbf{读出内容}。这是一种有启发性的解释，并不意味着每个真实神经元都恰好对应一条可分离的事实。
\end{conclusion}
